Table~1 lists the possible terms for various configurations of equivalent
$p$- and $d$-electrons. A number written beneath a term symbol gives the
number of terms of that type in the configuration, whenever there is more
than one. A configuration containing the largest possible number of
equivalent electrons ($s^2,p^6,d^{10},\ldots$) always has the term ${}^1S$.
Notice that the sets of terms are the same for two configurations when one
contains as many electrons as the other lacks for a filled shell. This follows
directly because a missing electron in a shell may be regarded as a
\emph{hole}, whose state is specified by the same quantum numbers as the state
of the absent electron.

\SourcePage{311}{314}
When Hund's rule is used to determine the normal term of an atom from its
known electron configuration, only the unfilled shell need be considered,
since the angular momenta of the electrons in filled shells cancel one
another. Suppose, for example, that an atom has four
